% ============================================
% 第3章：参考文献管理 + 图片上传 + 文档导出功能需求
% 负责人：C
% ============================================

\section{参考文献管理功能}

\subsection{功能描述}

参考文献管理模块为学生提供论文写作过程中参考文献条目的录入、编辑、检索和格式化输出能力。学生可以在系统内独立维护自己的参考文献库，在撰写论文正文时通过引用标注与文献条目关联，最终在导出文档时自动生成符合GB/T 7714-2015国家标准的参考文献列表。

\subsection{功能列表}

\begin{enumerate}
    \item \textbf{参考文献录入}：学生逐条录入参考文献信息，包括作者、标题、文献类型（期刊/书籍/会议/学位论文/网络资源）、出版年份、刊物名称、卷号、期号、页码、出版社、会议名称、学位授予单位、访问链接等字段；
    \item \textbf{文献编辑与删除}：支持对已录入的参考文献信息进行修改和删除操作；
    \item \textbf{文献检索}：支持按作者姓名和文献标题的关键词模糊搜索；
    \item \textbf{标准化格式化}：系统自动将文献信息按GB/T 7714-2015标准格式化为引用字符串，支持五种文献类型的格式化输出；
    \item \textbf{格式化列表预览}：学生可预览全部参考文献的格式化列表效果，确认格式正确再导出。
\end{enumerate}

\subsection{业务流程}

\begin{enumerate}
    \item 学生进入参考文献管理页面，系统展示已录入的文献列表（按年份降序排列）；
    \item 学生点击"新增文献"，弹出录入表单，填写作者、标题、文献类型及其他类型特有字段；
    \item 提交保存后，文献条目出现在列表中，系统同时生成并展示该条目的GB/T 7714格式化文本；
    \item 学生可随时编辑或删除已有文献条目；
    \item 导出论文时，系统自动将所有文献按格式化顺序列表写入文档末尾的"参考文献"章节。
\end{enumerate}

\subsection{数据描述}

参考文献数据以结构化方式存储在数据库中，每条文献包含以下核心信息：

\begin{itemize}
    \item \textbf{基本字段}：作者列表（多作者以逗号分隔）、文献标题、文献类型枚举（J/M/C/D/EB）、出版年份；
    \item \textbf{期刊专用（[J]）}：刊物名称、卷号、期号、页码；
    \item \textbf{书籍专用（[M]）}：出版社名称、出版地；
    \item \textbf{会议专用（[C]）}：会议名称；
    \item \textbf{学位论文专用（[D]）}：学位授予单位；
    \item \textbf{网络资源专用（[EB]）}：访问链接、引用日期。
\end{itemize}

\section{图片上传与管理功能}

\subsection{功能描述}

图片上传管理模块为学生提供论文写作过程中插入图片的功能支持。学生可从本地上传论文插图，系统对上传的图片进行格式校验和大小限制，存储后提供访问URL供前端在编辑器中展示，并在导出文档时将图片嵌入到对应的文档位置。

\subsection{功能列表}

\begin{enumerate}
    \item \textbf{图片上传}：学生通过文件选择器选择本地图片文件，上传至服务器；
    \item \textbf{格式校验}：系统自动校验上传文件的MIME类型，仅允许JPEG、PNG、GIF、WebP格式；
    \item \textbf{大小限制}：单张图片大小限制为10MB，超出限制时返回错误提示；
    \item \textbf{自动重命名}：上传后系统自动以UUID重命名文件，避免文件名冲突；
    \item \textbf{图片预览与访问}：上传成功后返回图片的访问URL，前端可直接嵌入编辑器或展示，支持在浏览器中直接查看图片文件；
    \item \textbf{图片删除}：学生可删除不再需要的已上传图片，同时清理服务器本地文件。
\end{enumerate}

\subsection{业务流程}

\begin{enumerate}
    \item 学生在编辑器中点击"插入图片"按钮，弹出文件选择对话框；
    \item 选择本地图片文件后，前端通过multipart/form-data格式POST请求上传至服务端；
    \item 服务端接收文件，校验MIME类型和文件大小，校验通过后以UUID重命名文件并保存至服务器本地磁盘；
    \item 服务端将图片元信息（原始名称、存储名称、文件路径、大小、类型）写入数据库，返回图片ID和访问URL；
    \item 前端使用返回的URL在编辑器中展示图片，用户保存后该URL嵌入论文章节的HTML内容中；
    \item 导出文档时，服务端根据URL读取本地图片文件，嵌入到DOCX或PDF的对应位置。
\end{enumerate}

\section{文档导出功能}

\subsection{功能描述}

文档导出模块是本系统的核心功能，负责将学生在Web编辑器中完成的论文内容，按照所选模板格式配置，生成标准DOCX（Word）和PDF格式文件供用户下载。导出过程为服务端异步操作，导出完成后提供文件下载链接。

\subsection{功能列表}

\begin{enumerate}
    \item \textbf{DOCX格式导出}：将论文标题、各章节段落、图片、表格、参考文献等内容，结合模板的字体、字号、行距、页边距等格式参数，通过Apache POI库生成标准.docx格式文件；
    \item \textbf{PDF格式导出}：在生成DOCX文件的基础上，通过LibreOffice无头转换模式将DOCX转换为PDF格式，保证排版一致性；
    \item \textbf{导出进度反馈}：导出任务执行期间，系统向前端反馈生成进度（生成中/完成/失败）；
    \item \textbf{文件下载}：导出完成后提供文件下载链接或直接流式传输文件。
\end{enumerate}

\subsection{导出内容范围}

导出的文档应完整包含以下内容：

\begin{enumerate}
    \item \textbf{封面}：论文标题、学生姓名、学院、专业、指导教师、日期等封面字段；
    \item \textbf{原创性声明}：预设的学术诚信声明文本；
    \item \textbf{中英文摘要}：论文的中文和英文摘要内容及关键词；
    \item \textbf{正文章节}：完整的多级章节结构，包含段落文字、图片、表格等素材；
    \item \textbf{参考文献}：按照GB/T 7714-2015标准格式排列的参考文献列表。
\end{enumerate}

\subsection{非功能性需求}

\begin{enumerate}
    \item \textbf{导出质量}：导出的DOCX文件使用宋体、Times New Roman等标准学术字体，中文内容无乱码或方框；
    \item \textbf{响应时间}：单次导出操作应在60秒内完成（论文规模不超过100页、图片10张以内）；
    \item \textbf{并发控制}：同一用户同一时间仅允许执行一次导出操作，防止重复提交；
    \item \textbf{错误处理}：导出失败时需明确告知失败原因（如图片文件缺失、磁盘空间不足等），并提供重试入口；
    \item \textbf{文件存储}：导出的文档文件按日期和用户ID组织存储，定期清理过期文件（保留7天）。
\end{enumerate}
